\section{鞅}

以下设$T\subseteq\mathbb{R}$。
\begin{definition}
	设$(X,\mathscr{F},P)$是概率空间。若$\{\mathscr{F}_t\}_{t\in T}$是一族$\mathscr{F}$的子$\sigma$域，且：
	\begin{equation*}
		s,t\in T,\;s\leqslant t\quad\Longrightarrow\quad\mathscr{F}_s\subseteq\mathscr{F}_t
	\end{equation*}
	则称$\{\mathscr{F}_t\}_{t\in T}$为该概率空间上的\gls{Filtration}或\textbf{信息流}。
\end{definition}
\begin{definition}
	设$\{f_t\}_{t\in T}$是概率空间$(X,\mathscr{F},P)$上取值于$(Y,\mathscr{A})$的随机过程，$\{\mathscr{F}_t\}_{t\in T}$是其所在概率空间上的滤过。若对每个$t\in T$，$f_t$都是从$(X,\mathscr{F}_t)$到$(Y,\mathscr{A})$的可测映射，则称该过程\textbf{适应}于此滤过。
\end{definition}
\begin{definition}
	设$\{f_t\}_{t\in T}$是概率空间$(X,\mathscr{F},P)$上取值于$(Y,\mathscr{A})$的随机过程。定义：
	\begin{equation*}
		\mathscr{F}_t\coloneq\sigma(f_s:s\in T,\;s\leqslant t)
	\end{equation*}
	称$\{\mathscr{F}_t\}_{t\in T}$为该过程的\gls{NaturalFiltration}。
\end{definition}
\begin{definition}
	设$\{f_t\}_{t\in T}$是概率空间$(X,\mathscr{F},P)$上取值于$(\mathbb{R},\mathcal{B}(\mathbb{R}))$的随机过程，$\{\mathscr{F}_t\}_{t\in T}$是该概率空间上的滤过。若：
	\begin{enumerate}
		\item $\{f_t\}_{t\in T}$适应于滤过$\{\mathscr{F}_t\}_{t\in T}$；
		\item 对任意的$t\in T$有$\operatorname{E}(|f_t|)<+\infty$;
		\item 对任意的$s,t\in T$且$s\leqslant t$有$\operatorname{E}(f_t\mid\mathscr{F}_s)=f_s\;$a.s.于$(X,\mathscr{F}_s,P)$；
	\end{enumerate}
	则称$\{f_t\}_{t\in T}$关于$\{\mathscr{F}_t\}_{t\in T}$为\gls{Martingale}。
\end{definition}